\chapter{entries/exits/daily/summary 怎么手算和检查}

\section{本章要解决的困惑}

fast backtest 跑完以后，不要只看一个总收益数字。一个可信的回测输出至少要能回答：

\begin{enumerate}[leftmargin=2em]
  \item 这次到底用了什么 profile？
  \item 哪些 entry 被触发，哪些被 admission 拒绝？
  \item target exposure 和 actual exposure 差多少？
  \item fixed60 是 mid 口径还是 executable 口径？
  \item 每天结果是否稳定，还是某一天撑起全部收益？
\end{enumerate}

\section{entries.parquet}

\code{entries.parquet} 是 entry 候选和 capacity 决策的明细表。重点字段：

\begin{longtable}{p{0.28\textwidth}p{0.62\textwidth}}
\toprule
字段 & 检查问题 \\
\midrule
\code{side} & 方向是否和 TFI 一致 \\
\code{cell} & R5/frames 分到哪个四象限 \\
\code{membership_set} & 哪些 trigger class 组合成 entry \\
\code{entry_bid/ask/mid} & 入场盘口是什么 \\
\code{entry_cross_bps} & 入场 cross 成本是多少 \\
\code{admission_accepted} & spread gate 是否允许 \\
\code{a_r5_raw} & entry 前闭合记忆是多少 \\
\code{base_weight/gamma} & 仓位倍率来自哪里 \\
\code{requested_exposure} & 策略想要多少 exposure \\
\code{actual_exposure} & 账本实际允许多少 \\
\code{clipped_exposure} & 因 3x 限制剪掉多少 \\
\bottomrule
\end{longtable}

真正有仓位的是：

\begin{lstlisting}
actual_exposure > 0
\end{lstlisting}

\section{exits.parquet}

\code{exits.parquet} 是退出明细。重点看：

\begin{longtable}{p{0.28\textwidth}p{0.62\textwidth}}
\toprule
字段 & 检查问题 \\
\midrule
\code{exit_profile} & \code{fixed60_mid}、\code{fixed60_taker} 还是别的 \\
\code{exit_reason} & 为什么退出 \\
\code{exit_bid/ask/mid} & 退出盘口是什么 \\
\code{raw_bps_mid_fixed60} & mid 口径固定 60 秒收益 \\
\code{net_bps} & 当前 profile 下净收益口径 \\
\code{actual_exposure} & 这笔退出对应多少仓位 \\
\bottomrule
\end{longtable}

如果 mid 很好但 executable 很差，不要说策略好。这可能只是 spread-cost mirage。

\section{daily.csv}

\code{daily.csv} 用来检查稳定性。至少看：

\begin{lstlisting}
entries
exits
actual_exposure
net_weighted_bps
worst_actual_leg_net_bps
max_open_exposure
admission_accepted_entries
admission_rejected_entries
\end{lstlisting}

如果三天总收益好，但其中一天很差，就不能只说“总收益不错”。当前系统样本本来就短，daily split 比总数更重要。

\section{summary.json}

\code{summary.json} 是 run 的封面。它必须告诉你：

\begin{lstlisting}
run_id
symbol
date range
policy
admission_profile
capacity_profile
exit_profile
actual_entries
exits
net_weighted_bps
timing_ms
\end{lstlisting}

如果两个 run 的 \code{exit_profile} 不同，一个是 \code{fixed60_mid}，一个是 \code{fixed60_taker}，不能直接比较收益。

\section{一个检查顺序}

建议每次按这个顺序：

\begin{lstlisting}
summary.json
-> profile_manifest.json
-> daily.csv
-> entries.parquet cell/admission/capacity
-> exits.parquet mid vs executable
-> worst day and worst leg
\end{lstlisting}

\section{手动检查点}

拿一笔 entry 和 exit，手算：

\[
Y^{exec,long}_{60}=10000\log\frac{exit\_bid}{entry\_ask},
\]

或：

\[
Y^{exec,short}_{60}=10000\log\frac{entry\_bid}{exit\_ask}.
\]

再乘以 \code{actual_exposure}，看是否能解释 daily 贡献。

\checkpoint{输出文件不是附属物。它们是你判断回测是否可信的证据层。}
